\section{What the positive route must now prove}
\label{sec:positive-route}

\subsection{Gate P: principal mass}

The exact obligation is
\[
 \mathfrak P_X
 =
 \sum_H\frac{2H}{q_X-1}M_H^{\mathrm{inv}}
 =
 \sum_{\substack{\beta\in\mathfrak B_X^{\mathrm{nc}}\\
                   1\le N_\beta\le q_X}}
 \frac{2H_{q_X}(N_\beta)}{q_X-1}
 \sum_{\substack{z\in I_\beta^\ast\\
                   q_X\nmid u_{\beta,z}}}
 w_{\beta,z}
 \le X^{o(1)}Q_X^2.
\]
For a branch of resolved \(z\)-fiber length \(N_\beta\ge3\),
\[
 \frac{H_{q_X}(N_\beta)}{q_X}
 \le\frac1{N_\beta}.
\]
This suggests opening the positive mass before grouping by \(H\)
and proving an inverse-length estimate.  It does \emph{not} justify
replacing \(M_H^{\mathrm{inv}}\) by an ambient count:
\(M_H^{\mathrm{inv}}\) contains the literal
row coefficients, content and projector masses, masks, and one
global normalization.

After the \(q_X\mid u_{\beta,z}\) atoms have been returned, the
remaining invertible endpoint atoms with \(N_\beta=1,2\) are purely
principal.  They receive no centered TPC-101 cost, so they must be
returned entirely through \(\mathfrak P_X\).  Calling them ``short''
does not make them negligible.

\subsection{Gate W: maximum opened-atom weight}

TPC-101 needs
\[
 W_X=\max w_C(j)\le X^{o(1)}
\]
to derive the diagonal gate from Gate P.  This must be checked by
writing one atom literally:
\begin{equation}
 w_C(j)
 =
 \left|
 \mathsf R_C(j)\,
 \mathsf C_C(j)\,
 \mathsf S_C(j)\,
 \mathsf N_X
 \right|.
\label{eq:literal-atom-template}
\end{equation}
Here \(\mathsf R_C\) contains the row coefficients,
\(\mathsf C_C\) the content and projector coefficient,
\(\mathsf S_C\) the local and smooth weights, and
\(\mathsf N_X\) the single global normalization.
An \(\ell^1\) bound for all rows does not by itself imply the desired
pointwise cap after a different normalization.  Conversely, if each
factor in \eqref{eq:literal-atom-template} has a subpower pointwise
envelope, Gate W is a local bookkeeping theorem rather than a new
M\"obius-cancellation theorem.

\subsection{Gate X: genuine cross-cell excess}

TPC-100 gives the exact collision support
\[
 A_Cm_Cj-A_Dm_Dk
 \equiv h_0(A_D-A_C)\pmod{q_X}.
\]
The remaining ledger is
\[
 \mathfrak X_X
 =
 \sum_Hg_{q_X}(H)
 \sqrt{|\mathfrak X_H^\circ|}.
\]
At the fully literal TPC-100 resolution, retaining
\((m,u,\sigma)\) already recovers
\(j=(\sigma u-h_0)/m\), so every fine cell is a singleton.
The resulting same-cell ``injectivity'' is therefore tautological
and supplies no dispersion gain.  All collisions between distinct
opened atoms---including atoms from the same parent \(\beta\)---stay
in the cross-cell term, and different cells may also induce the same
affine map.  A valid proof should therefore:
\begin{enumerate}[leftmargin=2em]
\item aggregate identical maps only after retaining a recoverable
  provenance bag and exact total weight;
\item separate identical-map mass from genuinely distinct-map
  collisions;
\item center with the physical baseline \(M_CM_D/(q_X-1)\); and
\item use the width factor \(g_{q_X}(H)\) before summing over \(H\).
\end{enumerate}

\begin{proposition}[Correct order of attack]
\label{prop:attack-order}
Within the positive route, Gate W should be audited before a new
diagonal-energy theorem, and Gate P before a global cross-cell
large sieve.
\end{proposition}

\begin{proof}
If Gates P and W hold, the diagonal term already follows from
\cref{thm:master}; a separate diagonal theorem would duplicate work.
If the actual principal term \(B_X\mathfrak P_X\) fails
polynomially, centered dispersion cannot repair the
\emph{termwise upper bound} in \cref{thm:master}, so a cross-cell
theorem alone would not close this three-ledger route.  A sharper
filter-sensitive positive theorem could still use cancellation
between the principal and nonprincipal pieces.  A lower bound for
\(\mathfrak P_X\) alone is not a stop certificate unless one also
knows \(B_X\ge X^{-o(1)}\).  The stated order tests the cheapest kill
conditions first.
\end{proof}

\begin{remark}[What counts as forward progress]
A proof of Gate W is \(\Lone\) bookkeeping progress.  A proof of
Gate P or Gate X for the actual growing census is closer to an
arithmetic estimate, but it becomes \(\Ltwo\) only after its exact
role in the full growing fixed-\(h_0\) coefficient and all remaining
reassembly losses is verified.
\end{remark}
